% ============================================================
%  Theory of Computation -- Homework Template
%  Compile with:  pdflatex hw-template.tex   (run twice)
% ============================================================
\documentclass[11pt]{article}

% ---------- Packages ----------
\usepackage[margin=1in]{geometry}
\usepackage{amsmath, amssymb, amsthm, mathtools}
\usepackage{enumitem}
\usepackage{booktabs}
\usepackage{array}
\usepackage{fancyhdr}
\usepackage{tikz}
\usetikzlibrary{automata, positioning, arrows.meta, calc}
\usepackage{algorithm}
\usepackage[noend]{algpseudocode}
\usepackage{xcolor}
\usepackage[hidelinks]{hyperref}

% ---------- Fill these in ----------
\newcommand{\studentname}{Your Name}
\newcommand{\coursename}{Theory of Computation}
\newcommand{\hwnumber}{1}
\newcommand{\duedate}{October 9, 2026}
\newcommand{\collaborators}{None}   % list anyone you worked with

% ---------- Header / footer ----------
\pagestyle{fancy}
\fancyhf{}
\lhead{\studentname}
\chead{\coursename}
\rhead{Homework \hwnumber}
\cfoot{\thepage}
\setlength{\headheight}{14pt}

% ---------- Problem / solution environments ----------
\newcounter{problem}
\newenvironment{problem}[1][]{%
  \refstepcounter{problem}\par\bigskip
  \noindent\textbf{Problem \theproblem}%
  \if\relax\detokenize{#1}\relax\else\ \textbf{(#1)}\fi
  \par\smallskip\noindent\itshape
}{\par\medskip}

\newenvironment{solution}{%
  \par\noindent\textbf{Solution.}\ \normalfont
}{\hfill$\blacksquare$\par\medskip}

% Sub-parts: \begin{parts} \item ... \end{parts}  -> (a), (b), ...
\newlist{parts}{enumerate}{1}
\setlist[parts]{label=(\alph*), leftmargin=*}

% Theorem-style environments
\theoremstyle{plain}
\newtheorem{theorem}{Theorem}
\newtheorem{lemma}[theorem]{Lemma}
\newtheorem{claim}[theorem]{Claim}
\newtheorem{corollary}[theorem]{Corollary}
\theoremstyle{definition}
\newtheorem{definition}[theorem]{Definition}

% ---------- Notation macros ----------
% Alphabets, strings
\newcommand{\Sig}{\Sigma}
\newcommand{\Gam}{\Gamma}
\newcommand{\blank}{\sqcup}                    % blank symbol
\newcommand{\eps}{\varepsilon}                 % empty string
\newcommand{\rev}[1]{#1^{R}}                   % reverse of a string
\newcommand{\chk}{\checkmark}                  % check mark on the tape

% Languages and machines
\newcommand{\lang}[1]{L(#1)}                   % L(M)
\newcommand{\comp}[1]{\overline{#1}}           % complement
\newcommand{\enc}[1]{\langle #1 \rangle}       % encoding <M>, <M,x>
\newcommand{\set}[2]{\left\{ #1 \,\middle|\, #2 \right\}}  % {x | ...}

% Named languages
\newcommand{\ATM}{A_{\mathrm{TM}}}
\newcommand{\HALT}{\mathrm{HALT}}
\newcommand{\EQTM}{\mathrm{EQ}_{\mathrm{TM}}}
\newcommand{\ETM}{E_{\mathrm{TM}}}

% Transitions:  \trans{q}{a}{p}{b}{R}  ->  delta(q,a) = (p,b,R)
\newcommand{\trans}[5]{\delta(#1, #2) = (#3, #4, #5)}
\newcommand{\qacc}{q_{\mathrm{accept}}}
\newcommand{\qrej}{q_{\mathrm{reject}}}

% Reductions
\newcommand{\mred}{\le_{m}}                    % mapping reducible
\newcommand{\red}{\le}                         % generic reduction

% Complexity / classes (for later in the course)
\newcommand{\cls}[1]{\mathsf{#1}}              % \cls{P}, \cls{NP}, \cls{RE}

% ---------- Tape diagram ----------
% \tape{a,b,c,\blank}{2}  draws cells a b c _ with the head over cell 2.
% Optional first arg labels the head with a state: \tape[q_1]{...}{2}
\newcommand{\tape}[3][]{%
  \begin{tikzpicture}[baseline=-0.5ex]
    \foreach \sym [count=\cellnum] in {#2} {
      \node[draw, minimum width=0.7cm, minimum height=0.7cm, inner sep=0pt]
        (c\cellnum) at (0.7*\cellnum, 0) {$\sym$};
      \xdef\lastcell{\cellnum}
    }
    \node[right=0pt of c\lastcell, minimum height=0.7cm] {$\cdots$};
    \draw[-{Latex}, thick] ($(c#3.north)+(0,0.55)$) -- (c#3.north);
    \node[above=0.55cm of c#3, inner sep=1pt] {$#1$};
  \end{tikzpicture}%
}

% ============================================================
\begin{document}

\begin{center}
  {\Large\bfseries \coursename\ --- Homework \hwnumber}\\[4pt]
  \studentname \quad\textbullet\quad Due \duedate\\[2pt]
  {\small Collaborators: \collaborators}
\end{center}
\hrule

% ============================================================
%  PROBLEM TEMPLATE 1: Designing a Turing machine
% ============================================================
\begin{problem}[Design a TM]
Specify a 1-tape Turing machine that accepts
\[
  L = \set{x\#y\#x}{x, y \in \{0,1\}^{+}}.
\]
\end{problem}

\begin{solution}
\textbf{Idea.} Check the format, then repeatedly check off the leftmost
unchecked symbol of the first block, remember it in the state, and check
off the matching symbol in the third block.

\textbf{High-level description.}
$M$ = ``On input $w$:
\begin{enumerate}[label=\arabic*.]
  \item Scan $w$; if it is not of the form $\{0,1\}^{+}\#\{0,1\}^{+}\#\{0,1\}^{+}$,
        \emph{reject}. Return the head to the left end.
  \item Move right past any $\chk$'s. If the head reads $\#$, go to stage 4.
  \item Remember the symbol $s \in \{0,1\}$ in the state and replace it
        with $\chk$. Move right past both $\#$'s and any $\chk$'s.
        If the symbol there is $s$, replace it with $\chk$ and return to the
        left end; go to stage 2. Otherwise (mismatch or $\blank$), \emph{reject}.
  \item Move right past $y$, the second $\#$, and any $\chk$'s.
        If the head reads $\blank$, \emph{accept}; otherwise \emph{reject}.''
\end{enumerate}

\textbf{Pseudocode} (optional).
\begin{algorithmic}[1]
  \State check format; \textbf{reject} if bad
  \Loop
    \State move right past all $\chk$
    \If{head reads $\#$} \textbf{break} \EndIf
    \State $s \gets$ symbol under head; write $\chk$
    \State move right past two $\#$'s and all $\chk$
    \If{symbol $= s$} write $\chk$ \Else\ \textbf{reject} \EndIf
    \State move left past two $\#$'s, then left to the last $\chk$
  \EndLoop
  \State skip $y$, $\#$, and $\chk$'s; \textbf{accept} iff head reads $\blank$
\end{algorithmic}

\textbf{Formal definition.} $M = (Q, \Sig, \Gam, \delta, q_0, \qacc, \qrej)$ with
$\Sig = \{0,1,\#\}$, $\Gam = \{0,1,\#,\chk,\blank\}$. Transitions not listed go to $\qrej$.

\begin{center}
\begin{tabular}{@{}lccccl@{}}
\toprule
State & Read & Write & Move & Next & Comment \\
\midrule
$q_0$       & $0,1$ & same   & R & $q_1$      & first block nonempty \\
$q_1$       & $0,1$ & same   & R & $q_1$      & \\
$q_1$       & $\#$  & $\#$   & R & $q_2$      & \\
$\vdots$    &       &        &   &            & (fill in remaining rows) \\
$c_0''$     & $0$   & $\chk$ & L & $q_{\text{back}}$ & match \\
$q_{\text{check}}$ & $\blank$ & $\blank$ & -- & $\qacc$ & \\
\bottomrule
\end{tabular}
\end{center}

\textbf{Example run} on $01\#1\#01$ (after first match):
\[
  \tape[q_{\text{back}}]{\chk,1,\#,1,\#,\chk,1,\blank}{5}
\]

\textbf{Correctness.} Each loop iteration checks off exactly one symbol of
the first block and the symbol in the same position of the third block,
rejecting on mismatch. The loop ends when the first block is fully checked;
stage 4 then accepts iff the third block has no unchecked symbols, i.e.\
iff both blocks have the same length and agree everywhere. Hence
$\lang{M} = L$.
\end{solution}

% ============================================================
%  PROBLEM TEMPLATE 2: Undecidability by reduction
% ============================================================
\begin{problem}[Undecidability]
Show that $\EQTM = \set{\enc{M_1, M_2}}{\lang{M_1} = \lang{M_2}}$ is undecidable.
\end{problem}

\begin{solution}
We reduce from $\ATM$, which is undecidable. Suppose for contradiction that
a TM $R$ decides $\EQTM$. Construct a TM $S$ deciding $\ATM$:

$S$ = ``On input $\enc{M, x}$:
\begin{enumerate}[label=\arabic*.]
  \item Construct $M_1$ = ``On input $y$: ignore $y$, simulate $M$ on $x$,
        and accept if $M$ accepts.''
  \item Construct $M_2$ = ``On input $y$: accept.''
  \item Run $R$ on $\enc{M_1, M_2}$ and output what $R$ outputs.''
\end{enumerate}

\textbf{Why it works.}
\begin{itemize}
  \item If $M$ accepts $x$, then $\lang{M_1} = \Sig^* = \lang{M_2}$, so $R$ accepts.
  \item If $M$ does not accept $x$, then $\lang{M_1} = \emptyset \ne \lang{M_2}$, so $R$ rejects.
\end{itemize}
$S$ always halts because $R$ does, so $S$ decides $\ATM$, a contradiction.
Therefore $\EQTM$ is undecidable.
\end{solution}

% ============================================================
%  PROBLEM TEMPLATE 3: Closure / recognizability proof
% ============================================================
\begin{problem}[Recognizability]
Prove: if $L$ and $\comp{L}$ are both recognizable, then $L$ is decidable.
\end{problem}

\begin{solution}
Let $M_1$ recognize $L$ and $M_2$ recognize $\comp{L}$. Define

$M$ = ``On input $w$: run $M_1$ and $M_2$ on $w$ in parallel, alternating
one step each. If $M_1$ accepts, \emph{accept}; if $M_2$ accepts, \emph{reject}.''

Every $w$ is in exactly one of $L$ or $\comp{L}$, so exactly one of $M_1, M_2$
eventually accepts. Thus $M$ always halts, and it accepts exactly $L$.
\end{solution}

% ============================================================
%  BLANK PROBLEM -- copy this block for new problems
% ============================================================
\begin{problem}
% Problem statement here.
\end{problem}

\begin{solution}
% Your answer here.
\end{solution}

\end{document}